\section{A baseline cost plus an amplitude penalty}
\label{sec:baseline-prox}
The statistical argument also permits a positive cost for producing any appreciable output. Its value may vary with the prior scale. The conclusion needed for the noise witness is unchanged: the posterior mean is close to a mixture of zero and a capped empirical predictor.

\begin{lemma}[A baseline-plus-amplitude posterior bound]
\label{lem:baseline-prox}
Let the input support be a finite set $\mathcal U$, let $|Y_i|\le2a$, and write $M(z)=\max_{x\in\mathcal U}|z(x)|$. Suppose $B\ge0$, $0<c\le4$, and $t>0$ satisfy the entropy bound
\[
 \mathfrak s\operatorname{KL}(\rho\|\mu)
 \ge B\Pr_\rho(M(z)>t)+c\,\mathbb E_\rho M(z)-e_0
\]
for every absolutely continuous comparison law, where $\mathfrak s>0$ is the entropy normalization and $s,\zeta,e_0,e_2\ge0$. Suppose the prior itself has expected empirical loss at most $\widehat R(0)+s$, and every function $v:\mathcal U\to[-2a,2a]$ has a comparison satisfying
\[
 \mathbb E\widehat R(z)\le\widehat R(v)+\zeta,
 \qquad \mathfrak s\operatorname{KL}(\rho_v\|\mu)
        \le B+4\|v\|_\infty+e_2.
\]
Let $\Pi$ be the Gibbs posterior minimizing
$\mathbb E\widehat R+\lambda\mathfrak s\operatorname{KL}$, and let $u_\lambda$ minimize
\[
 J_\lambda(u)=\widehat R(u)+\lambda c\|u\|_\infty
\]
on sample vectors constant on repeated inputs. Then $u_\lambda$ is a capped empirical-mean predictor. For some $p\in[0,1]$,
\[
 \|\mathbb E_\Pi z(S)-p u_\lambda\|_S
 \le t+\sqrt{4at+s+\zeta+
          \lambda\{2a(4-c)+e_0+e_2\}}.
\]
If $2a\le1$, the same bound holds for the mean of draws clipped to $[-1,1]$ and for clipping the raw posterior mean. The reference vector $p u_\lambda$ is nonpositive at every observed input with a nonpositive empirical label mean.
\end{lemma}
\begin{proof}
Set $A=\{M(z)>t\}$ and $p=\Pi(A)$. The low-output event has
\[
 \widehat R(z)\ge\widehat R(0)-2t\|y\|_S
               \ge\widehat R(0)-4at.
\]
The quadratic-gap inequality in Lemma~\ref{lem:canonical-prox} gives, on the high-output event,
\[
 \widehat R(z)+\lambda cM(z)
 \ge J_\lambda(u_\lambda)+\|z(S)-u_\lambda\|_S^2.
\]
Let $F$ be the posterior free energy. The entropy floor therefore implies
\begin{align*}
 F+\lambda e_0
 &\ge (1-p)\{\widehat R(0)-4at\}
       +p\{J_\lambda(u_\lambda)+\lambda B\}
       +\mathbb E_\Pi\mathbf1_A\|z(S)-u_\lambda\|_S^2.
\end{align*}
Extend $u_\lambda$ by zero at unobserved support points. Its sup norm is at most $2a$. The two variational comparisons give
\[
 F\le\min\left\{\widehat R(0)+s,
   J_\lambda(u_\lambda)+\lambda B+\zeta+
          \lambda[2a(4-c)+e_2]\right\}.
\]
For any two real numbers $A_0,A_1$ and nonnegative errors $r_0,r_1$,
\[
 \min(A_0+r_0,A_1+r_1)
 \le (1-p)A_0+pA_1+r_0+r_1.
\]
Subtracting the preceding inequalities bounds
$\mathbb E_\Pi\mathbf1_A\|z(S)-u_\lambda\|_S^2$ by the quantity under the square root. Meanwhile the low-output contribution to the mean has norm at most $t$. Cauchy--Schwarz controls the high-output contribution by that square root, proving the estimate.

Clipping contracts distance to $u_\lambda\in[-2a,2a]$ on the high event and cannot increase the low-event norm. Clipping the mean contracts distance to $p u_\lambda$. Finally the capped empirical-mean characterization from Lemma~\ref{lem:canonical-prox} proves the sign assertion.
\end{proof}

The baseline $B$ disappears from the approximation error. This cancellation is why a varying baseline cost does not itself protect a learner against the negative-singleton witness. It changes whether zero or a capped empirical predictor is favored, but both have the wrong sign on the same negative singleton observations when the true regression function is positive there.
